\section{Method Details and Mathematical Analysis}
\label{app:contracts}
\label{app:notation}

We first specify the interfaces that connect recollection, consequence assessment, and action generation, then establish how the shared gate controls recalled experience.
Training objectives and causal record construction complete the method specification; Appendix~\ref{app:symbols} indexes the core notation.

\subsection{Semantic Interfaces and Information Flow}
\label{app:effect_repr}

\paragraph{A common space for requirements and effects.}
The fixed descriptor map $\mathcal P_{\mathrm{eff}}$ defines a common space for historical effects $E_i$, requirements $r_t$, and predictions $\widehat E_i$ within each evaluated implementation.
Frozen DINOv2 features $\Phi$~\citep{oquab2024dinov2} are pooled to $2\times2$ cells of width $d_s$, with pre-state $S^-_{ij}$ and change $\Delta S_{ij}$.

\paragraph{Descriptor in the audited task specialists.}
\label{app:original_effect_projection}
For the audited Battery Try, Press Button, Cover Blocks, and Blocks Ranking Try specialists, the effect target is the uniform mean of the four recorded change tokens:
\[
 E_i=\mathcal P_{\mathrm{eff}}^{\mathrm{mean}}(S_i^-,\Delta S_i)
 =\frac14\sum_{j=1}^{4}\Delta S_{ij}.
\]
The demonstrated-query target $E_t^\star$ uses the same mean.
This vector is a coarse semantic change target; the retained spatial-token records and current visual/proprioceptive inputs remain separate from it.
Independent effects for a given checkpoint must be evaluated using its own descriptor map, without changing the map after observing outcomes.

\paragraph{Descriptor in the separate shared-stage variant.}
The separate shared-stage configuration uses the following change-weighted map:

\begin{equation}
 a_{ij}=\|\Delta S_{ij}\|_2^2/d_s,\quad
 \omega_{ij}=\frac{a_{ij}+\epsilon_p}{\sum_{k=1}^{4}a_{ik}+4\epsilon_p},\quad
 E_i=\sum_{j=1}^{4}\omega_{ij}(\Delta S_{ij}+0.05S^-_{ij}).
 \label{eq:effect_pooling}
\end{equation}
The stabilizer $\epsilon_p>0$ handles negligible changes, and $0.05S^-_{ij}$ preserves starting context.
Within this shared-stage variant, the same map gives $E_t^\star$; it differs from the uniform-mean target of the audited specialists above.

\paragraph{Connect current features to stored evidence.}
We read the same current frame at Video DiT layers 9 and 19, obtaining token streams $U_t^{(1)},U_t^{(2)}$ of width 3,072; these are two network depths, not two camera views.
Learned projections $P_1,P_2$ and mixing logits $b_1,b_2$ fuse them:
\begin{equation}
 W_t=\omega_1P_1U_t^{(1)}+\omega_2P_2U_t^{(2)},\qquad
 (\omega_1,\omega_2)=\operatorname{softmax}(b_1,b_2),
 \label{eq:world_fusion}
\end{equation}
Four learned Perceiver-style queries~\citep{jaegle2021perceiver} compress $W_t$ into $\widehat S_t\in\mathbb R^{4\times d_s}$, aligned to frozen $\Phi(o_t)$ during training.
This bridge connects the instruction- and state-conditioned video representation to stored evidence through the interfaces in Table~\ref{tab:app_visibility}.

\begin{table}[!htbp]
\centering\small
\caption{\textbf{Interfaces for consequence-aligned reuse.} Demonstrated actions and futures provide training targets; deployment reads use the inputs listed here.}
\label{tab:app_visibility}
\begin{tabularx}{\linewidth}{@{}>{\raggedright\arraybackslash}p{0.22\linewidth}>{\raggedright\arraybackslash}p{0.39\linewidth}>{\raggedright\arraybackslash}X@{}}
\toprule
Computation & Direct inputs & Output and role \\
\midrule
Requirement read & $W_t,\widehat S_t,q_t,\ell,\mathcal M_t$; retrieved groups masked & $R_t,r_t$: assessment reference \\
Action adaptation & $A_i,c_t,c_i^0,\tau_i$; current context and $R_t$ & $A_i^w,\mu_i$: aligned and adapted actions \\
Effect prediction & $\operatorname{stats}(\mu_i),E_i,\operatorname{pool}(W_t)$ & $\widehat E_i$: proposed effect \\
Ranking and gating & Candidate/context features, timing; recollection through the requirement; variant-specific gate inputs & $\ell_i^{\mathrm{rank}},\alpha_t,g_t$: selection and influence \\
Selected-event read & Requirement inputs plus selected pre-state/change tokens & $R_t^{\mathrm{pred}}$: transition context \\
Action-context readout & Adaptation states, stored change, predicted effect & $C_i$: projected action/effect context \\
Action DiT & $x_\sigma,\sigma,\mathcal C_t$ & $v_\theta$: velocity generating $A_t$ \\
\bottomrule
\end{tabularx}
\end{table}

\paragraph{Two reads, one Transition Adapter.}
The shared $G_{\mathrm{tr}}$ uses eight learned queries, two cross-attention/feed-forward blocks~\citep{vaswani2017attention}, type embeddings, validity masks, and the native robot-state token.
Masking retrieved groups gives $R_t\in\mathbb R^{8\times d_r}$ of adapter width $d_r$; $P_{\mathrm{req}}$ projects their mean to the candidate-independent reference $r_t\in\mathbb R^{d_s}$.
A second read adds the selected event's pre-state/change tokens to produce $R_t^{\mathrm{pred}}$.
We gate its increment over $R_t$ by Eq.~\eqref{eq:conditions} and apply affine $P_{\mathrm{ctx}}$ to reach decoder width.

\paragraph{From an action proposal to its predicted effect.}
The adapter $f_A$ combines aligned actions $A_i^w$, position and gripper timing $\tau_i$, current context, and $R_t$ to produce bounded $\mu_i$ in Eq.~\eqref{eq:adapter}.
Effect prediction uses mean command, terminal command, and mean absolute temporal variation:
\begin{equation}
 \operatorname{stats}(A)=\left[
 \frac1H\sum_{h=0}^{H-1}A_h,\ A_{H-1},\
 \frac1{H-1}\sum_{h=0}^{H-2}|A_{h+1}-A_h|
 \right]\in\mathbb R^{3d_a}.
 \label{eq:action_stats}
\end{equation}
For $H=1$, variation is zero.
World pooling applies a learned projection and validity-masked mean to $W_t$.
Together with $E_i$ and $\operatorname{stats}(\mu_i)$, it gives $\widehat E_i$ for the proposal preceding flow refinement.
Task context enters through $W_t$ and $R_t\rightarrow\mu_i\rightarrow\widehat E_i$.

Action-context tokens combine adaptation states, stored changes, and predicted effects.
Their affine projection to decoder width is included in $C_i$, so gating scales the entire contribution.


\paragraph{Compare direction and magnitude.}
For effect vectors $p,q$ above the cosine norm floor, Eq.~\eqref{eq:kappa} decomposes as
\begin{equation}
 1-\kappa(p,q)=\frac12\left\|\frac p{\|p\|_2}-\frac q{\|q\|_2}\right\|_2^2
 +\lambda_m\left|\log\frac{\|p\|_2+\epsilon}{\|q\|_2+\epsilon}\right|.
 \label{eq:comparison_geometry}
\end{equation}
The two terms measure angular mismatch and relative effect magnitude.
We use $\lambda_m=0.25$ and $\epsilon=10^{-6}$; safe cosine normalizes with denominator floor $\epsilon$, assigning zero cosine to a zero vector.
Our action summary compresses temporal order, and the descriptor compresses physical change; Appendix~\ref{app:effect_probes} tests whether this resolution preserves useful candidate ordering against independently executed outcomes.

\subsection{Recall Selection and Effective Gating}
\label{app:gate_details}

\paragraph{Retrieve, align, and select.}
For a unit-normalized frozen-DINO CLS vector $v_t\in\mathbb R^{768}$ and a one-hot task vector $e(\ell)\in\mathbb R^J$, the retrieval key is $k_t=2^{-1/2}[v_t;e(\ell)]$; stored keys use the same construction.
We apply cosine top-$K$ search to these keys, with $J=40,9,4$ for LIBERO, RMBench, and the physical task set, respectively.
The task block contributes similarity without excluding cross-task events.
Alignment $\mathcal W$ preserves relative EEF commands. For absolute joint commands, $c_t$ and $c_i^0$ denote current and stored starts in the same joint coordinates; their normalized difference shifts arm targets while preserving gripper commands.
When such matched start-state coordinates are unavailable, $\mathcal W$ leaves the stored command unchanged and the action adapter supplies the learned correction.
The adapter then applies the bounded residual in Eq.~\eqref{eq:adapter}.
This start-state alignment is expressed in the control coordinates; it does not establish object-frame or camera-frame geometric invariance.
Object displacement can change visual context, proposal suitability, and the correction needed for success even when the residual remains bounded.
The specialist reranker reads current/candidate context, adapted-action statistics, observed effects, and gripper timing to produce $\ell_i^{\mathrm{rank}}$.
In the shared-policy variant, projected episode history augments its query, and projected phase/ordinal features augment candidate context.
Both variants use recollection in the candidate-independent requirement and hence in consequence comparison.

\paragraph{Keep implementation and evidence aligned.}
\label{app:implementation_variants}
Our task-specialist benchmark models and the separately trained shared-policy control study implement the same requirement--proposal--influence interface.
Table~\ref{tab:implementation_variants} specifies their different encodings and selection/gate inputs.
We retain each checkpoint's own feature map and computation in the corresponding measurement; the additional shared-policy study does not replace the specialist benchmark results.

\begin{table}[!htbp]
\centering\small
\caption{\textbf{Implementation-specific interfaces.} Specialist interfaces correspond to RMBench task-specific training; shared-policy interfaces correspond to the update-3,000 Full/Pooled configuration.}
\label{tab:implementation_variants}
\begin{tabularx}{\linewidth}{@{}p{0.20\linewidth}XX@{}}\toprule
Interface & Task-specialist path & Shared-policy path\\\midrule
Effect descriptor & Uniform mean of four change tokens & Change-weighted descriptor in Eq.~\eqref{eq:effect_pooling}\\
Gripper timing & Four entries: first close/open phases and validity in the combined gripper stream & Four entries per gripper; eight in the bimanual configuration\\
Ranking context & Query/candidate context, adapted actions, observed effects, timing & Additionally projects episode history into the query and phase/ordinal into candidate context\\
Gate input & Selected score, score margin, compatibility, validity coordinate, deformation & Selected probability, probability margin, entropy, compatibility, validity coordinate, deformation, repetition\\
Event continuity & No additive event-order term at the specialist model boundary & Runtime supplies event-order prior and temporal source eligibility\\\bottomrule
\end{tabularx}
\end{table}

\Needspace{9\baselineskip}
\paragraph{Preserve useful event continuity in the shared-policy runtime.}
The following event-order extension supplies $b_{t,i}$ to the shared-policy selector; the audited specialist model uses rank plus consequence compatibility without this additive term.
Let the previously selected event come from demonstration $d_0$ at time $s_0$, and let $\Delta_i=s_i-s_0$ for candidates from that demonstration.
We define its recorded-order prior by
\[
 b_{t,i}=\begin{cases}
 0,&\text{no active event or a different demonstration},\\
 -\lambda_{\mathrm{reuse}},&\operatorname{demo}(i)=d_0,\ \Delta_i\leq0,\\
 \lambda_{\mathrm{forward}}\max(0,1-\Delta_i/F),&\operatorname{demo}(i)=d_0,\ \Delta_i>0,
 \end{cases}
\]
with nonnegative reuse penalty $\lambda_{\mathrm{reuse}}$, forward weight $\lambda_{\mathrm{forward}}$, and window $F>0$.
Same-or-earlier events are also source-ineligible; successor metadata may supply a continuation candidate.
This cursor tracks selected-event order, while recollection tracks actual interaction.
It resets each episode; advancement, null-step expiry, and continuation/truncation settings belong to the run specification in Table~\ref{tab:app_run_settings}.

\paragraph{Measure adaptation and repetition.}
Both variants measure adaptation by the selected residual. The additional repetition statistic below belongs to the shared-policy path.
For selected residual $\delta=\mu_{i^\star}-A_{i^\star}^w$, the smooth RMS deformation is
\[
 d=\sqrt{(Hd_a)^{-1}\|\delta\|_F^2+\epsilon_d^2}-\epsilon_d,
 \qquad \epsilon_d=10^{-6}.
\]
For $L$ executed-action summary slots, validity flags $m_j$, and stored repetition coordinates $r_j\in[0,1]$, we use
\[
 \zeta_t=\frac{\sum_{j=1}^La_jm_jr_j}{\max(1,\sum_{j=1}^La_jm_j)}.
\]
Recency weights $a_j$ are equally spaced from $0.5$ to $1$ ($0.5$ for one slot); empty history gives zero.
Each summary concatenates mean arm commands, terminal/net arm displacement, and terminal gripper values, with gripper entries zeroed in the arm blocks.
Relative commands use their sum for net displacement; absolute targets use offsets from the observed prefix-start state in both arm blocks.

We match a new executed-prefix summary to retained summaries by maximum cosine similarity, breaking ties by relative distance.
A repetition flag requires both similarity and distance thresholds to pass; it is zero without earlier history.
Observed progress clears it when frozen-feature change reaches the motion threshold, or both frozen-feature and VAE changes exceed the stationary threshold.
Reference defaults are $0.97$, $0.20$, $0.08$, and $0.03$; Table~\ref{tab:app_run_settings} records evaluated settings.

\paragraph{Turn selection evidence into shared influence.}
The specialist gate uses five inputs:
\begin{equation}
 \alpha_t^{\mathrm{sp}}=\operatorname{sigmoid} f_g^{\mathrm{sp}}\!
 \left([\ell_{t,i^\star}^{\mathrm{sel}},\Delta\ell_t,
 \kappa_{i^\star},\log(1+n),d]\right),\qquad
 g_t^{\mathrm{sp}}=\mathbf1[i^\star\in\mathcal I_t]\alpha_t^{\mathrm{sp}}.
 \label{eq:specialist_gate}
\end{equation}
Here $\Delta\ell_t$ is the selected score minus the strongest competing candidate score; with a single valid candidate it is zero.
The valid-candidate support is $n=1$, so this input is a validity coordinate.
The shared-policy variant instead forms probability-based confidence statistics and adds the repetition statistic $\zeta_t$:
for valid candidates, set $P_i=\operatorname{softmax}_i(\ell_{t,i}^{\mathrm{sel}}/\tau_{\mathrm{sel}})$ with $\tau_{\mathrm{sel}}>0$.
Let $P=P_{i^\star}$ and $\Delta P$ be its margin over the strongest competitor.
Normalized entropy is $\mathcal H=-\sum_iP_i\log P_i/\log|\mathcal I_t|$ for multiple candidates; a single candidate has zero entropy and an absent competitor of probability zero.
Its learned gate combines selection statistics, compatibility, validity coordinate $n$, deformation, and repetition:
\begin{equation}
 \alpha_t=\operatorname{sigmoid} f_g\!\left(
 [P,\Delta P,1-\mathcal H,\kappa_{i^\star},\log(1+n),d,\zeta_t]\right).
 \label{eq:rawgate}
\end{equation}
For each selected valid event, $n=1$ (and $n=0$ in padded slots), so $\log(1+n)$ is a fixed validity coordinate, not an estimated support count.
Equation~\eqref{eq:gate_loss} trains $\alpha_t$ as an imitation-utility score; eligibility and repetition determine its effective influence:
\begin{equation}
 g_t=\chi_t\,\mathbf1[\zeta_t<\zeta_{\max}]\,
 \alpha_t\exp(-\lambda_\zeta\zeta_t).
 \label{eq:effectivegate}
\end{equation}
For this shared-policy path, $\chi_t$ encodes candidate availability and eligibility: absent candidates, forced null, and forced negatives set it to zero; deployment also applies the raw-gate threshold and temporal eligibility.
The specialist path uses Eq.~\eqref{eq:specialist_gate}.
The coefficient $\lambda_\zeta\geq0$ sets repetition decay and $\zeta_{\max}\in[0,1]$ sets its cutoff.
In either path, the resulting $g_t$ controls both the stochastic source and candidate conditioning.

Empty retrieval bypasses argmax, zeros candidate payloads, and sets $g_t=0$.
The candidate-null interface retains the requirement context $P_{\mathrm{ctx}}(R_t)$, native context, visual K/V, and fixed conditioning masks, with zero action-context slots.
This common interface supports the exact withdrawal property below.

\subsection{Mathematical Properties of Selective Generative Reuse}
\label{app:theory}
\label{app:source_properties}

We establish exact withdrawal and bounded partial influence, then characterize the source and its flow-matching objective.
For withdrawal and source geometry, fix the current query and proposal $\mu=\mu_{i^\star}$, vectorize normalized actions in $\mathbb R^D$ with $D=Hd_a$, and write $\beta=1-\sigma_{\min}$ and $c(g)=1-\beta g$.
At prescribed gate $g\in[0,1]$, the source is $z_g=g\mu+c(g)\varepsilon_t$, where $\varepsilon_t\sim\mathcal N(0,I_D)$.

\subsubsection{Joint withdrawal and bounded action influence}
\label{app:influence_proof}

\paragraph{Proof of Proposition~\ref{prop:withdrawal}.}
Fix factual inputs, recollection, parameters, masks, solver schedule, and sampling randomness for one replan.
The candidate-independent read preserves $R_t$, native context, and visual K/V under candidate replacement.
At $g=0$, Eqs.~\eqref{eq:source} and~\eqref{eq:conditions} give $z_0=\varepsilon_t$, $R_t^{\mathrm{act}}=R_t$, and $C_t^{\mathrm{act}}=0$.
Every decoder input is identical; induction over solver steps therefore gives identical iterates and final actions.\hfill$\square$

\paragraph{Conditional episode equivalence to candidate-null generation.}
\label{app:candidate_null_episode_equivalence}
Consider two versions of the same trained policy, with identical parameters,
instruction, initial environment state, and initial retained factual state.
The first version retains the repertoire interface but fixes its effective
gate to zero at every replan. The second uses the candidate-null interface.
Assume matching decoder token layouts and masks, solver schedules, execution
prefixes, control and observation timing, stopping rules, and native
conditioning interfaces at each replan. Any wall-clock dependence of execution
or termination is held fixed by the comparison protocol. Couple all sampling, environment,
observation, and controller randomness between the two versions.
Assume further that every retained state reaching an ungated decoder input
is updated by the same candidate-independent function of factual observations
and executed commands. Candidate-dependent bookkeeping may differ only when
it has no route to an ungated decoder input.

At the first replan, joint withdrawal gives identical Gaussian sources,
recollection-grounded requirements, zero candidate conditioning, and matched
native and visual context. The fixed solver therefore produces identical
action chunks and executed prefixes. Identical executed commands under the
coupled environment and controller randomness yield identical next states
and observations. The assumed factual-state update then yields identical
retained inputs for the following replan. Induction proves pathwise equality
of the complete executed trajectories. Consequently every measurable success
criterion on those trajectories has equal probability under the two versions.

This is conditional equivalence to the candidate-null version of the same
trained policy, not an empirical comparison with a separately trained
retrieval-free policy. It does not assert that a learned gate becomes zero
when an entire skill is absent from a nonempty bank. That detection property
requires a bank-wide skill-absence protocol and is separate from the
zero-gate premise used here.


\paragraph{Partial admission within a fixed proposal branch.}
Hold $\mu,R_t,R_t^{\mathrm{pred}},C_{i^\star}$, factual context, and masks fixed, and vary $g$ with a shared noise draw.
For the linear part $P_{\mathrm{ctx}}^{\mathrm{lin}}$ of the affine context projection, define
\begin{equation}
 h_t=\left[\operatorname{vec}\!\left(P_{\mathrm{ctx}}^{\mathrm{lin}}
 (R_t^{\mathrm{pred}}-R_t)\right);\operatorname{vec}(C_{i^\star})\right].
 \label{eq:context_direction}
\end{equation}
Gating the entire projected candidate contribution gives $\mathcal C_g-\mathcal C_0=gh_t$ in the varying coordinates.
Assume continuous-flow solutions exist on the unit interval and the following bound holds uniformly over their trajectories and contexts for finite $L_x,L_C\geq0$:
\begin{equation}
 \|v_\theta(x,\sigma;\mathcal C)-v_\theta(y,\sigma;\mathcal C')\|
 \leq L_x\|x-y\|+L_C\|\mathcal C-\mathcal C'\|,
 \label{eq:lipschitz_assumption}
\end{equation}
With $\psi(L)=(e^L-1)/L$ for $L>0$ and $\psi(0)=1$, we obtain
\begin{equation}
 \|A_t(g)-A_t(0)\|\leq g\left[
 e^{L_x}\|\mu-\beta\varepsilon_t\|+L_C\psi(L_x)\|h_t\|\right]
 \equiv gK_t.
 \label{eq:action_influence_bound}
\end{equation}
Here $A_t(g)$ is the continuous normalized output before discontinuous command processing, and $K_t$ is independent of $g$ on the fixed branch.

\noindent\emph{Proof.}
Set $s=1-\sigma$, $Y_g(s)=x_{1-s}(g)$, and $D(s)=\|Y_g(s)-Y_0(s)\|$.
Then $\mathrm dY_g/\mathrm ds=-v_\theta(Y_g,1-s;\mathcal C_g)$, and $D(0)=g\|\mu-\beta\varepsilon_t\|$.
Equation~\eqref{eq:lipschitz_assumption} yields
\[
 D(s)\leq D(0)+\int_0^s[L_xD(u)+gL_C\|h_t\|]\,\mathrm du.
\]
Gronwall's inequality gives
$D(s)\leq e^{L_xs}D(0)+gL_C\|h_t\|\int_0^se^{L_x(s-u)}\mathrm du$;
at $s=1$ this is Eq.~\eqref{eq:action_influence_bound}, including $L_x=0$.\hfill$\square$

For a shared Euler grid with positive steps $\Delta_k$ summing to one, the recurrence
$D_{k+1}\leq(1+L_x\Delta_k)D_k+gL_C\Delta_k\|h_t\|$
gives the same bound using $\prod_k(1+L_x\Delta_k)\leq e^{L_x}$, under the same regularity on the iterates.
Thus the shared gate controls both source and conditioning through a linearly shrinking influence envelope.
The comparison holds within fixed candidate, mask, and eligibility choices, before physical execution; output differences inside the envelope may vary nonmonotonically.
Appendix~\ref{app:gate_probes} checks the corresponding zero-gate withdrawal in the implementation.

\subsubsection{Geometry of the experience-informed source}
\label{app:source_geometry}

Writing $\varepsilon=\varepsilon_t$, the source law conditioned on a proposal and gate fixed before noise is drawn is
\begin{equation}
 \nu_g=\operatorname{Law}(z_g\mid\mu,g)=\mathcal N(g\mu,c(g)^2I_D),\qquad
 c(g)\geq\sigma_{\min}>0.
 \label{eq:source_law}
\end{equation}
The positive floor keeps the source nondegenerate at full acceptance; across queries, these Gaussian sources can form a mixture.
For a fixed demonstrated action $A^\star$, the zero-mean cross term gives
\begin{equation}
 \mathbb E_\varepsilon\|z_g-A^\star\|_F^2
 =\|g\mu-A^\star\|_F^2+c(g)^2D,
 \label{eq:geometry}
\end{equation}
Writing this expected distance as $J(g)$, its change from a Gaussian start is
\begin{equation}
 J(g)-J(0)=g^2(\|\mu\|_F^2+\beta^2D)
 -2g(\langle\mu,A^\star\rangle_F+\beta D).
 \label{eq:source_distance_difference}
\end{equation}
The expression separates the proposal shift from noise contraction, clarifying the two contributions of the shared source.
For the candidate-null law $\nu_0=\mathcal N(0,I_D)$, the quadratic Wasserstein distance is
\begin{equation}
 W_2(\nu_g,\nu_0)=g\sqrt{\|\mu\|_F^2+\beta^2D}.
 \label{eq:source_wasserstein}
\end{equation}
The common-noise coupling attains squared cost $g^2\|\mu\|_F^2+(c(g)-1)^2D$; Cauchy--Schwarz bounds every coupling's centered cross moment by $Dc(g)$, establishing optimality.
Hence source displacement scales linearly with admission, while the learned flow determines the resulting action.

\subsubsection{Conditional flow matching with experience-dependent sources}
\label{app:flow_proof}

Our specialization of conditional flow matching~\citep{lipman2023flowmatching,liu2023rectifiedflow} models the joint distribution of demonstrated actions and retrospective sources under the declared training corruptions.

\begin{warmproposition}[Flow matching with experience-dependent sources]
\label{prop:conditional_flow}
Fix the retrospective maps and consider full-valid action chunks.
Let $(A^\star,Z,\mathcal C)$ have any joint distribution with finite second moments, where $\mathcal C$ is precisely the decoder-visible condition.
Define $X_\sigma=(1-\sigma)A^\star+\sigma Z$ and $V=Z-A^\star$.
Assume $\mathbb E[w(\sigma)\|V\|^2]<\infty$. For independently sampled $\sigma$ with positive density and $w(\sigma)>0$, the unrestricted minimizer of the weighted squared flow loss is, almost everywhere,
\begin{equation}
 u_\sigma(x;c)=\mathbb E[V\mid X_\sigma=x,\mathcal C=c].
 \label{eq:optimal_conditional_field}
\end{equation}
The conditional path law $p_\sigma(\cdot\mid c)$ satisfies the continuity equation with field $u_\sigma$ in the weak sense.
If its ODE and continuity equation are well posed backward on each $[\delta,1]$, $\delta>0$, sampling from $p_1(\cdot\mid c)$ and integrating the field recovers $p_\delta(\cdot\mid c)$, with weak endpoint limit $p_0(\cdot\mid c)$.
\end{warmproposition}

\noindent\emph{Proof.}
For fixed $\sigma,c$, conditional expectation gives
\begin{equation}
 \mathbb E[\|v(X_\sigma,\sigma;c)-V\|^2\mid c]
 =\mathbb E[\|v(X_\sigma,\sigma;c)-u_\sigma(X_\sigma;c)\|^2\mid c]
 +\mathbb E[\|V-u_\sigma(X_\sigma;c)\|^2\mid c].
 \label{eq:conditional_orthogonality}
\end{equation}
The cross term vanishes because $\mathbb E[V-u_\sigma(X_\sigma;c)\mid X_\sigma,c]=0$.
Positive scheduler weights and fixed action-dimension normalization preserve this pointwise minimizer.
For a smooth compactly supported test function $\varphi$,
\begin{equation}
 \frac{\mathrm d}{\mathrm d\sigma}\mathbb E[\varphi(X_\sigma)\mid c]
 =\mathbb E[\nabla\varphi(X_\sigma)^\top V\mid c]
 =\mathbb E[\nabla\varphi(X_\sigma)^\top u_\sigma(X_\sigma;c)\mid c].
 \label{eq:weak_continuity}
\end{equation}
Finite second moments and bounded $\nabla\varphi$ justify differentiation under the expectation.
This is the weak continuity equation $\partial_\sigma p_\sigma+\nabla\!\cdot(p_\sigma u_\sigma)=0$.
Uniqueness on $[\delta,1]$ identifies the ODE law initialized from $p_1$ with $p_\sigma$; mean-square convergence $X_\delta\to A^\star$ gives the weak endpoint limit.\hfill$\square$

This result connects experience-conditioned initialization to the population objective under the decoder's actual information interface.
Transport assumes an optimal field and matching source sampling; deployment eligibility, finite data, joint updates, model approximation, and discretization are assessed through the empirical policy evaluations.

\subsection{Auxiliary Supervision and Gradient Flow}
\label{app:learning_details}

Our targets connect each retrospective decision to observable evidence: recorded effects supervise semantic prediction, demonstration agreement supervises ranking and adaptation, and selected-proposal agreement supervises raw utility.

\paragraph{Masks and distances.}
Let $m_{t,h}\in\{0,1\}$ mark valid action steps and $f_t\in\{0,1\}$ mark availability of the future observation.
For queries with at least one valid step, define
\begin{equation}
 d_A(U,V)=\frac{\sum_{h=0}^{H-1}m_{t,h}\|U_h-V_h\|_2^2}{d_a\sum_{h=0}^{H-1}m_{t,h}},
 \qquad d_E(p,q)=\frac1{d_s}\|p-q\|_2^2.
 \label{eq:distances}
\end{equation}
Candidate sums and softmaxes use $\mathcal I_t$.
We write $\langle\cdot\rangle_{\mathcal S}$ for a query-set mean, with empty means equal to zero.
Sets $\mathcal A,\mathcal F,\mathcal V$ contain action-valid queries, those with future targets, and those with valid candidates, respectively; $\mathcal F,\mathcal V\subseteq\mathcal A$.
Set $\mathcal N$ contains forced hard negatives.
Missing futures retain available action supervision.

\paragraph{Requirement and semantic bridge.}
For prediction $p$ and detached target $y$, our semantic loss combines direction and log-magnitude alignment:
\begin{equation}
 \ell_{\mathrm{sem}}(p,y)=1-\cos(p,y)+\lambda_m\operatorname{SmoothL1}
 \!\left(\log(1+\|p\|_2),\log(1+\|y\|_2)\right).
 \label{eq:semantic_loss}
\end{equation}
Here $\operatorname{SmoothL1}$ denotes smooth absolute error.
We set $\mathcal L_{\mathrm{req}}=\langle\ell_{\mathrm{sem}}(r_t,\operatorname{sg}(E_t^\star))\rangle_{\mathcal F}$ and average $1-\cos(\widehat S_{t,j},\operatorname{sg}(\Phi(o_t)_j))$ over queries and four cells for $\mathcal L_{\mathrm{bridge}}$.
The frozen encoder supplies stable observed targets. Training uses a smooth robust penalty on $\log(1+\|\cdot\|_2)$; selection uses the absolute difference of stabilized log magnitudes in Eq.~\eqref{eq:kappa}. The coefficient $\lambda_m$ weights the comparison of predicted and required effect scales.

\paragraph{Ranking and action adaptation.}
For $t\in\mathcal V$, aligned stored actions and historical effects define detached utility weights:
\begin{equation}
 u_i=-d_A(A_i^w,A_t^\star)-\lambda_E f_t d_E(E_i,E_t^\star),\qquad
 y_i=\operatorname{softmax}_{i\in\mathcal I_t}(u_i/\tau_u).
 \label{eq:utility}
\end{equation}
Here $\tau_u>0$ sets concentration and $\lambda_E\geq0$ balances effect agreement, omitted when $f_t=0$.
These labels depend on deterministic alignment and recorded evidence, remaining fixed as the learned proposal and prediction change.
We optimize
\begin{equation}
 \begin{aligned}
 \mathcal L_{\mathrm{rank}}&=\left\langle D_{\mathrm{KL}}\!\left(
 \operatorname{sg}(y)\,\|\,\operatorname{softmax}(\ell^{\mathrm{rank}})\right)\right\rangle_{\mathcal V},\\
 \mathcal L_{\mathrm{adapt}}&=\left\langle\sum_{i\in\mathcal I_t}
 \operatorname{sg}(y_i)d_A(\mu_i,A_t^\star)\right\rangle_{\mathcal V\setminus\mathcal N}.
 \end{aligned}
 \label{eq:ranking_adaptation}
\end{equation}
The Kullback--Leibler loss $D_{\mathrm{KL}}$ trains relative preference among valid candidates; the weighted action loss trains bounded adaptation, excluding forced negatives.

\paragraph{Effect-supervision variants and their scope.}
\label{app:effect_target_scope}
The audited task-specialist runs for Battery Try, Press Button, Cover Blocks, and Blocks Ranking Try use event-factual effect supervision.
For each valid retrieved event, the effect branch is trained against that event's own recorded pre/post-transition descriptor $E_i$, rather than assigning the current demonstration's effect $E_t^\star$ to every candidate.
The current demonstration still supervises requirement prediction and contributes to ranking and utility targets. The descriptor for these runs is the uniform change mean in Appendix~\ref{app:original_effect_projection}.
These targets must be distinguished from the independently observed effect of executing an alternative action in the current query state.
In particular, because the historical descriptor is an input to the effect branch, reconstructing $E_i$ can minimize event-factual supervision without establishing transfer to the current state.
Thus event-factual targets alone do not identify current-state counterfactual consequences.

\paragraph{Shared-stage factual term and candidate regularizer.}
The separate shared-stage configuration uses the following alternative effect objective, distinguishing its shared-query regularizer from the event-specific targets of the four specialists above.
An observed transition $(o_t,A_t^\star,o_{t+H})$ trains the shared head with its historical-effect input set to zero:
\begin{equation}
 \mathcal L_{\mathrm{fact}}=\left\langle
 \ell_{\mathrm{sem}}\!\left(f_E(\operatorname{stats}(A_t^\star),0,
 \operatorname{pool}(W_t)),\operatorname{sg}(E_t^\star)\right)
 \right\rangle_{\mathcal F}.
 \label{eq:factual_effect_loss}
\end{equation}
For $\mathcal R=(\mathcal V\cap\mathcal F)\setminus\mathcal N$, we also train the candidate-input path with detached utility weights:
\begin{equation}
 \mathcal L_{\mathrm{cand}}=\left\langle\sum_{i\in\mathcal I_t}
 \operatorname{sg}(y_i)\ell_{\mathrm{sem}}(\widehat E_i,\operatorname{sg}(E_t^\star))
 \right\rangle_{\mathcal R}.
 \label{eq:candidate_regularizer}
\end{equation}
We use $\mathcal L_{\mathrm{effect}}=(\mathcal L_{\mathrm{fact}}+\mathcal L_{\mathrm{cand}})/2$ when $\mathcal R$ is nonempty, and $\mathcal L_{\mathrm{fact}}$ otherwise.
The candidate term is a shared-query regularizer: predicting $E_t^\star$ for every candidate minimizes it alone for nonzero targets above the norm floor.
Utility weights change each candidate's influence, not the shared target, and therefore do not provide the observed consequence of executing each alternative.
The factual term anchors the shared head on demonstrated transitions; it does not by itself rule out candidate-insensitive predictions on retrieved actions.
Appendix~\ref{app:effect_probes} therefore evaluates candidate sensitivity against independently executed effects, comparing $\widehat E_i$ with $r_t$ and $E_i$.

\paragraph{Utility targets and retrieval corruptions.}
Selected-proposal agreement defines the detached soft target for raw utility:
\begin{equation}
 y_t^g=\operatorname{sg}\!\left[\exp\!\left(-
 \frac{d_A(\mu_{i^\star},A_t^\star)+\lambda_g f_t d_E(\widehat E_{i^\star},E_t^\star)}{\tau_g}
 \right)\right],\qquad
 \mathcal L_{\mathrm{gate}}=\left\langle\operatorname{BCE}(\alpha_t,y_t^g)\right\rangle_{\mathcal A}.
 \label{eq:gate_loss}
\end{equation}
Here $\tau_g>0$ sets sharpness and $\lambda_g\geq0$ weights effect agreement, omitted when $f_t=0$.
Null and forced-negative targets are zero; binary cross-entropy (BCE) includes all action-valid rows and is evaluated from logits for stability.
The target measures imitation utility and evolves as the proposal and effect models are updated.

For this shared-stage variant, training mixes normal retrieval, candidate dropout, forced null, and hard negatives selected by demonstrated action/effect disagreement.
Forced negatives train ranking and rejection with $g_t=0$ on both generation inputs.
For these negatives, adaptation and candidate-effect losses are masked; factual transitions continue to supervise requirement and effect prediction.

\paragraph{Flow objectives and gradient routes.}
For action flow, we draw $u\sim\mathcal U(0,1)$ and set $\sigma=5u/(1+4u)$.
The scheduler weight $w(\sigma)$ is proportional to $\exp[-2(\sigma-1/2)^2]$ after subtracting its minimum on the 1,000-point training grid and normalizing the grid mean to one.
At inference, Euler integration uses the same shifted grid at $u_k=1-k/N$ for $k=0,\ldots,N$, with $N=10$ except $N=8$ for Observe and Pick Up (Table~\ref{tab:app_run_settings}).
Action-flow training applies this weight with the valid-action mask; a separate video-only pass computes $\mathcal L_{\mathrm{FM}}^v$ from observed-video targets with its corresponding schedule.
The action pass obtains visual K/V from current-frame processing.
We detach $z_t$ in both $x_\sigma$ and $v^\star$, and detach $g_t$ in conditioning while preserving gradients through the token representations.
Thus action-flow gradients train permitted conditioning paths, auxiliary losses train retrospective decisions, and gate supervision also reaches its feature inputs.
Retrieval and argmax choose discrete indices; selected continuous readouts remain trainable.

We update Action DiT, retrospective modules, the native state projection, and rank-16 bottleneck adapters at video layers 9 and 19.
The VAE, text encoder, $\Phi$, and remaining video-backbone parameters stay fixed.
Table~\ref{tab:app_run_settings} specifies the evaluated loss weights and training settings.

\subsection{Event Construction and Causal Recollection}
\label{app:memory_details}
\label{app:execution}

\paragraph{Build intact action--effect exemplars.}
Each repertoire event contains $H$ recorded commands and their boundary observations from a training demonstration.
Window starts combine regular anchors, a tail-aligned window, interaction-change peaks, and independently retained gripper transitions.
We normalize RMS changes in actions, robot state, and available visual features by their median and median absolute deviation (MAD), average the signals, and select local maxima with nearby-peak suppression.
Only complete windows at the recorded control rate are retained, storing $(S_i^-,A_i,\Delta S_i)$, start state, timing, and provenance separately from retrieval keys.
The bank stays fixed at evaluation; offline retrieval excludes the query episode and content-identical visual features.

\paragraph{Retain observed interaction.}
Recollection starts with the initial observed anchor.
After each executed prefix, we form four change components $c_{t,k}$: any gripper transition, commanded-direction change, frozen-feature change, and VAE-latent change.
Feature changes use
$d_{\mathrm{rel}}(u,v)=\min\{1,\|u-v\|_2/(\|u\|_2+\|v\|_2+\epsilon_0)\}$ with $\epsilon_0=10^{-8}$ and zero for two negligible inputs.
Direction change is $(1-\cos)/2$ between consecutive terminal/net arm displacements; the first prefix averages adjacent nonzero command-vector changes.
Absolute-joint summaries use start-relative offsets, and relative controls use displacements.
With nonnegative weights of positive total weight, $s_t^{\mathrm{write}}=\sum_kw_kc_{t,k}/\sum_kw_k$.

We write an event when this score exceeds
$\operatorname{median}(\mathcal Q)+1.4826\lambda_{\mathrm{MAD}}\operatorname{MAD}(\mathcal Q)$,
clipped to configured bounds; warm-up uses the lower bound.
The bounded queue $\mathcal Q$ contains earlier scores, and receives the current score only after the decision.
This order makes offline prefix reconstruction and online retention use identical available evidence.

\paragraph{Compress and read bounded recollection.}
New events have unit mass and post-observation representatives.
We keep the initial anchor separately and protect the latest event.
At capacity, we merge the most similar eligible adjacent interior pair by cosine similarity of flattened features, choosing the earliest pair on ties.
For representatives $f_a,f_b$ of masses $m_a,m_b$, the merged feature is $(m_af_a+m_bf_b)/(m_a+m_b)$ with combined mass, older pre-state, and newer post-state.
Repertoire actions remain intact.
Readout uses role/relative-age embeddings for initial, historical, latest, and executed-action records, plus separate validity bits for each gripper's opening and closing times.
Table~\ref{tab:app_run_settings} records capacities, weights, and thresholds.

\paragraph{Read, commit, and execute causally.}
At replan $t$, we read earlier records and already-executed commands alongside current $o_t,q_t$.
After inference, we commit the current observation with its preceding commands; newly generated commands enter history after execution, and the next observation supplies their feedback.
We reuse frozen retrieval features and reset recollection each episode.
Identical observed-and-executed prefixes therefore produce identical records under fixed settings: induction preserves the write evidence, earlier-score queue, and eligible merge pairs at every update.

\subsection{Core Notation and Module Index}
\label{app:symbols}

Table~\ref{tab:notation_index} prioritizes the main method's recurring symbols and essential module interfaces.
Candidate quantities refer to the current query $t$ unless indicated; hats mark learned representations, stars mark demonstration targets, and $i^\star$ is the selected index.
Local proof and diagnostic variables are defined where used; numerical settings appear in Table~\ref{tab:app_run_settings}.

\begingroup\small
\setlength{\tabcolsep}{4pt}
\setlength{\LTcapwidth}{\linewidth}
\renewcommand{\arraystretch}{1.08}
\begin{longtable}{@{}>{\raggedright\arraybackslash}p{0.265\linewidth}>{\raggedright\arraybackslash}p{0.53\linewidth}>{\raggedright\arraybackslash}p{0.15\linewidth}@{}}
\caption{\textbf{Core notation by computational role.} Definition references connect the main method to its detailed specification.}\label{tab:notation_index}\\
\toprule
Symbol or module & Meaning and role & Definition \\
\midrule
\endfirsthead
\multicolumn{3}{@{}l}{Table~\thetable{} (continued)}\\
\toprule
Symbol or module & Meaning and role & Definition \\
\midrule
\endhead
\multicolumn{3}{@{}l}{\textbf{Interaction, experience, and actions}}\\
$t,o_t,q_t,\ell$ & Replan time, RGB observation, robot state, instruction. & \S\ref{sec:method} \\
$\mathcal M_t,\mathcal B$ & Current-episode recollection; fixed cross-episode repertoire. & \S\ref{sec:memory} \\
$H,h_{\mathrm{exec}},d_a$ & Predicted horizon, executed prefix, per-command action width. & \S\ref{sec:method} \\
$e_i,S_i^-,\Delta S_i$ & Stored event, pre-state features, observed feature change. & Eq.~\eqref{eq:bank} \\
$A_i,A_i^w,\mu_i$ & Stored, aligned, and adapted action chunks in $\mathbb R^{H\times d_a}$. & Eq.~\eqref{eq:adapter} \\
$A_t,A_t^\star$ & Generated action chunk and demonstrated target. & \S\ref{sec:reuse}, \ref{sec:learning} \\
\midrule
\multicolumn{3}{@{}l}{\textbf{Semantic representations and learned interfaces}}\\
$\Phi,d_s,\mathcal P_{\mathrm{eff}}$ & Frozen visual encoder, semantic width, fixed effect descriptor map. & Appendix~\ref{app:effect_repr} \\
$W_t,\widehat S_t$ & Fused world tokens; learned semantic tokens aligned to $\Phi(o_t)$. & App.~\ref{app:effect_repr} \\
$G_{\mathrm{tr}},P_{\mathrm{req}},R_t,r_t$ & Shared adapter and requirement projection; output tokens and effect-space forecast. & Eq.~\eqref{eq:required} \\
$c_t,c_i^0,\mathcal W,f_A,\rho,s$ & Compatible action starts, alignment, residual adapter, bound, and scales. & Eq.~\eqref{eq:adapter} \\
$f_E,\widehat E_i$ & Effect head and its prediction for adapted proposal $\mu_i$. & Eq.~\eqref{eq:effect} \\
$E_i,E_t^\star,E_{t,i}^{\mathrm{obs}}$ & Historical effect, query-demonstration target, and independently executed candidate effect. & Eqs.~\eqref{eq:effect_target},\eqref{eq:effect_probe_target} \\
\midrule
\multicolumn{3}{@{}l}{\textbf{Selection and shared control of generation}}\\
$\mathcal I_t,K,i^\star$ & Valid candidates, retrieval count, selected index. & \S\ref{sec:compatibility} \\
$\kappa_i,\lambda_m,\lambda_c$ & Effect compatibility, magnitude weight, compatibility weight in selection. & Eqs.~\eqref{eq:kappa},\eqref{eq:selection} \\
$\ell_i^{\mathrm{rank}},b_{t,i},\ell_{t,i}^{\mathrm{sel}}$ & Learned relevance, optional shared-policy event-order prior, combined score. & Eq.~\eqref{eq:selection} \\
$f_g,\alpha_t,g_t$ & Utility head, raw utility, effective gate shared by both generation inputs. & App.~\ref{app:gate_details} \\
$z_t,\varepsilon_t,\sigma_{\min}$ & Stochastic source, standard Gaussian draw, retained-noise floor. & Eq.~\eqref{eq:source} \\
$R_t^{\mathrm{pred}},C_i$ & Selected-event transition read and projected action/effect context. & \S\ref{sec:reuse} \\
$R_t^{\mathrm{act}},C_t^{\mathrm{act}}$ & Gated transition and action-context inputs. & Eq.~\eqref{eq:conditions} \\
$P_{\mathrm{ctx}},\mathcal C_t$ & Affine context projection; full decoder-visible condition, including masks. & App.~\ref{app:effect_repr} \\
$v_\theta,x_\sigma,\sigma$ & Action DiT velocity, intermediate action, flow time from source ($1$) to action ($0$). & \S\ref{sec:reuse} \\
\midrule
\multicolumn{3}{@{}l}{\textbf{Learning and supervision}}\\
$d_A,d_E,\operatorname{sg}$ & Normalized action/effect squared distances; stop-gradient. & Eq.~\eqref{eq:distances}; \S\ref{sec:reuse} \\
$\mathcal L_{\mathrm{FM}}^a,\mathcal L_{\mathrm{FM}}^v$ & Action-flow and video-flow objectives. & Eqs.~\eqref{eq:action_fm},\eqref{eq:loss} \\
$\mathcal L_j,\lambda_j,\lambda_v$ & Six auxiliary losses and auxiliary/video weights in the joint objective. & Eq.~\eqref{eq:loss} \\
\bottomrule
\end{longtable}
\endgroup
